\chapter{KẾT LUẬN VÀ KIẾN NGHỊ}
\label{ch:conclusion}

\textit{(Nội dung chi tiết của chương này sẽ được phát triển đầy đủ trong giai đoạn tiếp theo):}

\section{Tổng kết các kết quả đạt được}
\begin{itemize}
    \item Đã xây dựng và khảo sát toàn diện mô hình toán học Lotka--Volterra tổng quát 2 loài tích hợp yếu tố khai thác và kinh tế sinh thái cho hệ cá mòi dầu -- cá vược sọc tại vịnh Chesapeake.
    \item Đã chứng minh chặt chẽ tính đặt đúng của mô hình (tính dương và bị chặn toàn cục trong miền $\Omega$).
    \item Đã chỉ ra sự khác biệt bản chất và những hiểm họa tiềm ẩn khi áp dụng chính sách MSY đơn loài cho cá mồi; chứng minh sự cần thiết của việc giảm 50\% hạn ngạch đánh bắt cá mòi dầu để bảo vệ kẻ săn mồi.
    \item Đã phân tích tính khả thi và giới hạn của chính sách MEY và RMY đối với hệ thống đa loài.
    \item Đã định lượng mức độ nhạy cảm của các tham số then chốt thông qua phương pháp vi phân trực tiếp.
\end{itemize}

\section{Hàm ý chính sách quản lý cho vịnh Chesapeake}
\begin{itemize}
    \item Chuyển đổi căn bản từ quản lý đơn loài sang Quản lý Nghề cá Dựa trên Hệ sinh thái (EBFM).
    \item Thiết lập các Điểm tham chiếu sinh thái (ERP) nghiêm ngặt đối với cá mòi dầu Đại Tây Dương để bảo tồn nguồn thức ăn cho cá vược sọc và các loài thủy sản thượng tầng.
    \item Kiểm soát chặt chẽ đội tàu lưới vây công nghiệp và hạn ngạch của các tập đoàn đánh bắt lớn.
\end{itemize}

\section{Đánh giá hạn chế và hướng phát triển tiếp theo}
\begin{itemize}
    \item Mô hình hiện tại chưa xét đến cấu trúc tuổi (\textit{age-structured models}) và kích thước cá thể.
    \item Cần mở rộng mô hình sang mạng lưới nhiều hơn 2 loài (kết hợp cá xanh, cá yếu, sinh vật phù du và các loài chim di cư).
    \item Tích hợp các yếu tố môi trường ngẫu nhiên (\textit{stochastic differential equations}) và tác động của biến đổi khí hậu trong các nghiên cứu mở rộng tiếp theo.
\end{itemize}
